\documentclass[UserManual.tex]{subfiles}
\begin{document}
\begin{sloppypar}
\section{Markov-Chain Monte Carlo Generation of the Posterior}\label{sec:mcmc}

\subsection{Overview}
The \textit{SMOOTH\_HOME}{\tt /software/MainPrograms/} directory contains the source code for the main program, {\tt smoothy\_mcmc\_main.cc}. When compiled it creates an executable, \textit{SMOOTH\_HOME}{\tt /software/bin/smoothy\_mcmc}, which then performs an MCMC analysis. This program will invoke the emulator in its exploration of the model-parameter space. It will write the trace to file and perform and record some analysis of the trace the User might find useful. Even if the User chooses a different MCMC approach then the one provided with the {\it Smooth Emulator} software, viewing the source code might prove useful in demonstrating how to invoke the emulator from within a C++ code. 

\subsection{Experimental Measurement Information}
Once the emulator is tuned and before it is applied to a Markov Chain investigation of the likelihood, the software needs to know the experimental measurement and uncertainty. That information must be entered in the {\tt smooth\_data/Info/experimental\_info.txt} file. The file should have the format:
\vspace*{-8pt}{\tt \begin{Verbatim}[commandchars=\\\{\}]
 observable_name_1  value_1    exp_uncertainty_1   theory_uncertainty_1
 observable_name_2  value_2    exp_uncertainty_2   theory_uncertainty_2
 observable_name_3  value_3    exp_uncertainty_3   theory_uncertainty_3
 observable_name_4  value_4    exp_uncertainty_4   theory_uncertainty_4
 .
\end{Verbatim}
}\vspace*{-8pt}
The second column provides the measured value and the third gives the experimentally reported uncertainty. The fourth column lists the uncertainty inherent to the theory, due to missing physics or to some shortcoming of the model. For example, even if a model has all the parameters set to the exact value, e.g. some parameter of the standard model, the full-model can't be expected to exactly reproduce a correct measurement given that some physics is likely missing from the full model. This should not include any random point-by-point uncertainty of the theory -- that is already included through the {\tt ALPHA} option in the emulator uncertainty. For the MCMC software, the relevant uncertainty incorporates both of these, plus the emulator uncertainty. Only the combination of all three, added in quadrature, affects the outcome. We emphasize that this last file is not needed to train and tune the emulator. It is only needed once one performs the MCMC search of parameter space. 

\subsection{Accounting for Uncertainties}

The log-likelihood, $LL$, for the MCMC generation is assumed to be of a simple form. Summing over the observables $I$,
\begin{eqnarray*}
\sigma_{I,{\rm tot}}^{2}&=&\sigma_{I,{\rm exp}}^2+\sigma_{I,{\rm theory}}^2+\sigma_{I,{\rm emulator}}^2,\\
LL&=&-\sum_{I}\frac{(Y_{I,{\rm exp}}-Y_{I,{\rm emu}})^{2}}{2\sigma_{I,{\rm tot}}^{2}}-\ln(\sigma_{I,{\rm tot}}).
\end{eqnarray*}
The likelihood behaves as $e^{LL}$. This form presumes that the uncertainties are Gaussian and that the uncertainties for the various observables are independent. The last term above derives from the normalization factor for the Gaussian. One can view \href{smoothdraft.pdf}{\tt smoothdraft.pdf} for more details.

Correlated errors are not explicitly addressed in {\it Smooth Emulator} software. However, one can always add additional model parameters that describe contributions to multiple observables due to some unknown factors, such as an experimental efficiency or some neglected physics. This can be a convenient and sound perspective from which one can incorporate both experimental and theoretical systematic uncertainties. For example, if all the masses measured by some experiment have a contribution to their uncertainties that either moves all of them up or down together, that contribution can be treated as an unknown model parameter, $\delta m$. This model parameter is simply added to all the other masses. From the Bayesian perspective, the model is trying to describe what the experiment will measure, not the ``true'' value, hence adding $\delta m$ is part of the model. Sometimes such parameters are called ``nuisance'' parameters, because they usually describe effects that are not of direct interest to the reader, but do impact the degree to which one might constrain the parameters of greater interest. However, the moniker ``nuisance'' can be rather unfortunate  because although some parameter might be a nuisance to one scientist, e.g. the efficiency of the detector, that same parameter might represent the life's work and passion of another scientist, who might find the constraint of such a parameter of great interest. Nuisance and non-nuisance parameters are treated on equal footings.

\subsection{MCMC Options}

Next, one needs to edit the options file {\tt\$\{MY\_ANALYSIS\}/smooth\_data/Options/mcmc\_options.txt}. An example file is:
\vspace*{-8pt}{\tt
\begin{Verbatim}
 # This is for the MCMC search of parameter space
 # (not for the emulator tuning)
 #LogFileName smoothlog.txt # comment out for interactive running
 # If success rate approaches zero, reduce MCMC_METROPOLIS_STEPSIZE, if it approaches 1.0, increase it
 MCMC_METROPOLIS_STEPSIZE 0.06
 MCMC_IGNORE_EMULATOR_ERROR false
 MCMC_RANSEED  12345
 MCMC_NBURN 5
 MCMC_NSKIP 10000
 MCMC_NTRACE 100000
\end{Verbatim}
}\vspace*{-8pt}
As was the case with {\it Smooth Emulator}, each option has a default value. The {\tt CMCMC} class has a parameter map object. This map stores the options, and uses them to set the corresponding data in the {\tt CMCMC} object. The values of {\tt MCMC\_NBURN}, {\tt MCMC\_NTRACE} and {\tt MCMC\_NSKIP} are stored by the {\tt CMCMC} parameter map, but do not correspond to any members of the class. When the User wishes to set the number of trace points in the main program, they do so by changing the arguments to the {\tt PerformTrace(..)} function. The User can set these arguments through the parameter map functionality, which can provide a convenient way to set the arguments through the Options file. See Chapter \ref{sec:mcmcsource} below for an example of how this is used.The options are described below.
\subsubsection{{\tt LogFileName (if commented out will write to screen)}}
To run interactively, leave this line commented out. Otherwise, output will be directed to the designated file.
\subsubsection{{\tt MCMC\_METROPOLIS\_STEPSIZE (default 0.05)}}
Metropolis algorithms require taking random steps in $\vec{\theta}$ space. If the steps are small, it takes longer to explore the space, but if the steps are very long, the success rate of the Metropolis steps becomes low. Maximum efficiency occurs when the Metropolis success rate, which is provided during running, is near 50\%. Rates of 20\% or 80\% are also fine, but if the rate is only a few percent, the User should reduce the parameter, and if the success rate becomes close to 100\%, the User should increase the parameter. If the acceptance rate is outside this range, the MCMC procedure will still provide a correct representation of the poster, but is might require more points to do so.
\subsubsection{{\tt MCMC\_IGNORE\_EMULATOR\_ERROR (default false)}}
If this flag is set to {\tt true} the emulator uncertainty will be ignored. This significantly increases the speed of the MCMC procedure, but should not be done if the emulator error is significant.
\subsubsection{{\tt MCMC\_RANSEED (default 12345)}}
Sets seed for random number generator.
\subsubsection{{\tt MCMC\_NBURN, MCMC\_NTRACE, MCMC\_NSKIP (no defaults)}}
These are different than the other options in that they do not correspond to any member of the {\tt CMCMC} class. Instead, they are simply stored by the {\tt CparameterMap} object of the {\tt CMCMC} class for convenience. The variables describe the number of random steps to be taken during the burn-in stage (where the trace is not being recorded), the number of steps in the trace to be recorded once the burn-in stage has ended, and the number of steps between recorded steps. Because successive steps are highly correlated it makes sense to not record every step. For example, if the Metropolis success rate is 50\%, then half the points are identical to the previous point.

The relevant variables are set when the CMCMC object calls the function {\tt PerformTrace(int Ntrace,int Nskip)} in the main program. When writing a main program, the User can set {\tt Ntrace} or {\tt Nskip} through the parameter map, which is a member of the {\tt CMCMC class}. This is illustrated below in Chapter \ref{sec:mcmcsource}, where the source code for the main MCMC program is described. The User can always simply call the {\tt PerformTrace} function with whatever arguments they wish and never use the options stored in the parameter map. The parameter map functionality is described in Sec. \ref{subsec:parameterMap}.

\subsection{Running the MCMC Program}\label{subsec:mcmcrunning}
To run the provided MCMC program, move to the analysis directory, and run the program {\tt mcmc}. Output should look something like this:
\vspace*{-8pt}{\tt
\begin{Verbatim}[commandchars=\\\{\}]
\textit{MY_ANALYSIS}% \textcolor{darkred}{\textit{SMOOTH_HOME}/bin/smoothy_mcmc}
At beginning of Trace, LL=-68.764478
At end of trace, best LL=1.563806
Best Theta=
0.249554  0.153237  0.190531  0.210907  0.058929  0.230998  
Metropolis success percentage=54.090000
finished burn in
At beginning of Trace, LL=-5.477040
finished 10%
finished 20%
finished 30%
finished 40%
finished 50%
finished 60%
finished 70%
finished 80%
finished 90%
finished 100%
At end of trace, best LL=1.678685
Best Theta=
0.269108  0.099867  0.185094  0.204939  0.037417  0.207398  
Metropolis success percentage=53.609600
writing, ntrace = 100001
writing, ntrace = 100001
\end{Verbatim}
}\vspace*{-8pt}
For the sake of numerical efficiency {\tt ``Metropolis success percentage''} should be in the range of 50\%. If the efficiency is very near 100\%, it suggests the step sizes might be too small to most efficiently explore the entire model-parameter space. If the efficiency is near zero, the step size might be too large and too few successful Metropolis steps will be successful. This affects only the efficiency, not the validity, so any success percentage between 10\% and 90\% should suffice. In the output, {\tt ``LL''} refers to the log-likelihood. At the end of the burn-in, one hopes that best value of {\tt LL} is not much lower than the best value found from the entire trace. If not, one should probably increase the value of {\tt MCMC\_NBURN}. The values of {\tt ``Best Theta''} refer to the point in the trace with the highest {\tt LL}. 

Information about the trace is found in the files located in the directory {\tt smooth\_data/mcmc\_trace/}:
\begin{itemize}
\item {\tt trace\_theta.txt}: A list of the scaled model-parameter values, $\vec{\theta}$, from the posterior sampling.
\item {\tt trace\_X.txt}: A list of the non-scaled model-parameter values, $\vec{X}$, from the posterior sampling.
\item {\tt xbar\_thetabar.txt}: The average of the parameter values ($\vec{X}$), and the scaled values ($\vec{\theta}$) from the posterior.
\item {\tt CovThetaTheta.txt}: This gives the $N_{\rm par}\times N_{\rm par}$ covariance matrix $\langle\delta \theta_i\delta\theta_j\langle$, describing the size and shape of the points in the trace.
\item {\tt CovThetaTheta\_eigenvalues.txt}: That eigenvalues of that matrix
\item {\tt CovThetaTheta\_eigenvecs.txt}: The eigenvectors
\item {\tt ResolvingPower.txt}: An $N_{\rm pars}\times N_{\rm obs}$ matrix describing the influence of each observable in resolving each model parameter.
\end{itemize}

\subsection{Reviewing the MCMC Source Code}\label{sec:mcmcsource}

Finally, we review the source code in \textit{SMOOTH\_HOME}{\tt /software/MainPrograms/mcmc\_main.cc}:
{\tt
\begin{Verbatim}[commandchars=\\\{\}]
int main()\{
   NBandSmooth::CSmoothMaster master;
   master.TuneAllY();
   NBandSmooth::CMCMC mcmc(&master);
   unsigned int Nburn=mcmc.parmap->getI("MCMC_NBURN",1000);  // Steps for burn in
   unsigned int Ntrace=mcmc.parmap->getI("MCMC_NTRACE",1000); // Record this many points
   unsigned int Nskip=mcmc.parmap->getI("MCMC_NSKIP",5); // Only record every Nskip^th point
      
   mcmc.PerformTrace(1,Nburn);   
   CLog::Info("finished burn in{\textbackslash}n");
   
   mcmc.PruneTrace(); // Throws away all but last point
   mcmc.PerformTrace(Ntrace,Nskip);
   mcmc.WriteTrace(); // Writes trace
   mcmc.EvaluateTrace();

   return 0;
\}
\end{Verbatim}
}

This is mostly self-explanatory. If one wishes to avoid writing out the trace, the line {\tt mcmc.WriteTrace} can be deleted. If the User wishes to run the code in batch mode, the output can be directed to a file, {\tt mcmc\_log.txt}, by adding the line\\
\hspace*{20pt}{\tt LogFileName mcmc\_log.txt} \\
to the options file {\tt smooth\_data/Options/mcmc\_options.txt}. The line {\tt mcmc.EvaluateTrace()} will evaluate the trace and calculate the resolving power and covariances.

If one has previously trained the emulator, one can use the known value of $\Lambda$ rather than recalculating it. Normally, this saves very little time, but might be noticeable when the number of model parameters is large.

\subsection{The MCMC header file}

In the main repository one can find a header file, \textit{SMOOTH\_HOME}{\tt /software/inlude/msu\_smooth/mcmc.h}. Perusing this file is a good way to understand the provided functionality of the CMCMC class. 
{\tt
\begin{Verbatim}[commandchars=\\\{\}]
class CMCMC\{
   public:
      bool IGNORE_EMULATOR_ERROR;
      CparameterMap *parmap;
      CSmoothMaster *master;
      Crandy *randy;
      
      CMCMC();
      CMCMC(CSmoothMaster *master);
      unsigned int NPars,NObs;
      vector<vector<double>> trace;
      string trace_filename,Xtrace_filename;
      double stepsize;
      
      void ClearTrace(); // erases trace info so one can start over, resets at theta=0.
      void PruneTrace(); // erases trace, except for last point
      
      void PerformTrace(unsigned int Ntrace,unsigned int Nskip);
      void PerformMetropolisTrace(unsigned int Ntrace,unsigned int Nskip);
      void WriteTrace();
      void ReadTrace();
      void EvaluateTrace();
      
      Eigen::VectorXcd stepvec,stepvecprime,dTdTEigenVals;
      Eigen::MatrixXd dThetadTheta;
      Eigen::MatrixXcd dTdTEigenVecs;

      void CalcLL(vector<double> &theta,double &LL);
      CLLCalcSmooth *llcalc;
      static CPriorInfo *priorinfo;
 \};
\end{Verbatim}
}
The functionalities that are most likely to be incorporated are the functions {\tt ClearTrace(), PruneTrace(), PerformTrace(), WriteTrace(), EvaluateTrace()} and {\tt CalcLL()}. These functions are mostly self-explanatory. The function {\tt WriteTrace()} will write the two files with trace information described in Sec. \ref{subsec:mcmcrunning}. The function {\tt CalcLL(vector<double> \&theta,double \&LL)} calculates the posterior log-likelihood, {\tt LL}, for the point $\vec{\theta}$.

The function {\tt EvaluateTrace()} considers the trace, assuming it has been calculated. After evaluating the trace it finds and writes the averages and covariances mentioned in Sec. \ref{subsec:mcmcrunning}. It also writes the file describing the resolving power for each observable w.r.t. the constraint of each observable. All the trace-evaluation output files are written to files in the director{\tt smooth\_data/MCMC/} directory.

\end{sloppypar}

\end{document}